% =============================================================
% ANEXO I. Relación del TFG con los ODS de la Agenda 2030
% Obligatorio según la plantilla de la EPSG.
% =============================================================

\chapter*{ANEXO I. Relación del TFG/TFM con los ODS de la agenda 2030}
\addcontentsline{toc}{chapter}{ANEXO I. Relación del TFG/TFM con los ODS de la agenda 2030}

Grado de relación del trabajo con los Objetivos de Desarrollo Sostenible (ODS).

\renewcommand{\arraystretch}{1.35}
\begin{longtable}{|p{6.7cm}|>{\centering\arraybackslash}p{1.5cm}|>{\centering\arraybackslash}p{1.5cm}|>{\centering\arraybackslash}p{1.5cm}|>{\centering\arraybackslash}p{1.8cm}|}
\hline
\textbf{Objetivos de Desarrollo Sostenibles} & \textbf{Alto} & \textbf{Medio} & \textbf{Bajo} & \textbf{No procede} \\
\hline
\endfirsthead
\hline
\textbf{Objetivos de Desarrollo Sostenibles} & \textbf{Alto} & \textbf{Medio} & \textbf{Bajo} & \textbf{No procede} \\
\hline
\endhead
ODS 1.\ Fin de la pobreza. & & & & X \\
\hline
ODS 2.\ Hambre cero. & & & & X \\
\hline
ODS 3.\ Salud y bienestar. & & & & X \\
\hline
ODS 4.\ Educación de calidad. & & & X & \\
\hline
ODS 5.\ Igualdad de género. & & & & X \\
\hline
ODS 6.\ Agua limpia y saneamiento. & & & & X \\
\hline
ODS 7.\ Energía asequible y no contaminante. & & & & X \\
\hline
ODS 8.\ Trabajo decente y crecimiento económico. & & & & X \\
\hline
ODS 9.\ Industria, innovación e infraestructuras. & & X & & \\
\hline
ODS 10.\ Reducción de las desigualdades. & & & & X \\
\hline
ODS 11.\ Ciudades y comunidades sostenibles. & & & & X \\
\hline
ODS 12.\ Producción y consumo responsables. & & & X & \\
\hline
ODS 13.\ Acción por el clima. & & & & X \\
\hline
ODS 14.\ Vida submarina. & & & & X \\
\hline
ODS 15.\ Vida de ecosistemas terrestres. & & & & X \\
\hline
ODS 16.\ Paz, justicia e instituciones sólidas. & & & & X \\
\hline
ODS 17.\ Alianzas para lograr objetivos. & & & & X \\
\hline
\end{longtable}

\section*{Justificación}
\addcontentsline{toc}{section}{Justificación}

Este Trabajo de Fin de Grado es un proyecto de investigación en el área de la informática, sin una intervención directa de carácter social, ambiental o económico. Su relación con los Objetivos de Desarrollo Sostenible es, por tanto, de naturaleza metodológica y se concentra en unos pocos objetivos.

La relación más clara es con el \textbf{ODS 9 (Industria, innovación e infraestructuras)}, valorada como media. El trabajo contribuye a la capacidad de investigación e innovación tecnológica: adapta una metaheurística clásica, el algoritmo genético, a un problema de búsqueda combinatoria con restricciones duras y evaluación costosa, y desarrolla para ello soluciones que no son específicas del dominio. Los operadores que trabajan a nivel de pack y el sistema de torneo que mantiene acotado el coste de la evaluación son aportaciones transferibles a otros problemas de búsqueda con restricciones estructurales y función objetivo cara. La meta 9.5 del ODS 9, que persigue reforzar la investigación científica y la capacidad tecnológica, es la que mejor recoge esta contribución.

La relación con el \textbf{ODS 4 (Educación de calidad)} se valora como baja. Como resultado de un trabajo académico, la memoria documenta el diseño, la calibración y el análisis de un algoritmo evolutivo de principio a fin, incluidos los caminos que no funcionaron y el motivo de su descarte. Ese recorrido completo, con sus aciertos y sus límites, tiene valor didáctico como material de apoyo para el estudio de las metaheurísticas aplicadas.

La relación con el \textbf{ODS 12 (Producción y consumo responsables)} también se valora como baja. La evaluación por simulación tiene un coste computacional alto, y el trabajo dedica parte de su esfuerzo a contenerlo. El cambio del sistema de torneo round-robin al sistema Swiss redujo en más de un 90\% el número de partidas simuladas por generación, y las líneas de trabajo futuro plantean un modelo sustituto que reduciría todavía más esa carga. Esa preocupación por un uso eficiente de los recursos de cómputo conecta, de forma modesta, con un consumo responsable.

El resto de objetivos no guarda una relación directa con la naturaleza del trabajo, por lo que se han marcado como no procedentes.
